% ─────────────────────────────────────────────────────────────────────────────
% Requiere en el preámbulo:
%   \usepackage{hanging}   % para \begin{hangparas}
%   \usepackage{hyperref}  % para \url
% ─────────────────────────────────────────────────────────────────────────────

\section{Referencias}

El desarrollo de esta tarea se apoyó en conocimiento previo de la materia, consulta de bibliografía clásica de algoritmos y conversaciones con modelos
de inteligencia artificial, a los que se proveyó como contexto el código fuente y los notebooks de medición.
Las herramientas de IA fueron utilizadas exclusivamente como asistentes de programación en los entornos Visual Studio Code y Windsurf; el análisis,
la interpretación de resultados y la redacción son propios.

% ── Bibliografía académica ────────────────────────────────────────────────────
\subsection{Bibliografía}

\begin{hangparas}{1.3cm}{1}

Cormen, T. H., Leiserson, C. E., Rivest, R. L., y Stein, C. (2022).
\textit{Introduction to algorithms} (4.ª ed.). MIT Press.

Hardy, G. H., y Wright, E. M. (2008).
\textit{An introduction to the theory of numbers} (6.ª ed.). Oxford University Press.

Knuth, D. E. (1997).
\textit{The art of computer programming: Vol. 2. Seminumerical algorithms}
(3.ª ed.). Addison-Wesley.

Sedgewick, R., y Wayne, K. (2011).
\textit{Algorithms} (4.ª ed.). Addison-Wesley Professional.

\end{hangparas}

% ── Modelos de inteligencia artificial ───────────────────────────────────────
\subsection{Herramientas de inteligencia artificial}

Los siguientes modelos de lenguaje fueron empleados como asistentes de
programación durante el desarrollo de la tarea, siguiendo las recomendaciones
de citación de herramientas de IA de la séptima edición de las normas APA
(American Psychological Association, 2020).

\begin{hangparas}{1.3cm}{1}

Anthropic. (2025). \textit{Claude Haiku 4.5} [modelo de lenguaje de gran escala].
\url{https://www.anthropic.com/news/claude-haiku-4-5}

Anthropic. (2025). \textit{Claude Sonnet 4.6} [modelo de lenguaje de gran escala].
\url{https://www.anthropic.com/news/claude-sonnet-4-6}

Cognition AI. (2025). \textit{SWE-1.5} [modelo de lenguaje de gran escala].
\url{https://cognition.ai/blog/swe-1-5}

\end{hangparas}

% ── Fundamentos técnicos ──────────────────────────────────────────────────────
\subsection{Fundamentos técnicos aplicados}

\begin{hangparas}{1.3cm}{1}

Cormen, T. H., Leiserson, C. E., Rivest, R. L., y Stein, C. (2022).
Memoización y programación dinámica. En
\textit{Introduction to algorithms} (4.ª ed., pp. 362--380). MIT Press.

Knuth, D. E. (1997).
Diccionarios y tablas de dispersión. En
\textit{The art of computer programming: Vol. 3. Sorting and searching}
(2.ª ed., pp. 506--549). Addison-Wesley.

Sedgewick, R., y Wayne, K. (2011).
Criba de Eratóstenes y variantes. En
\textit{Algorithms} (4.ª ed., pp. 204--206). Addison-Wesley Professional.

Umbach, J. (2021). \textit{LaTeX Workshop} (versión 10.x) [extensión para
Visual Studio Code]. GitHub.
\url{https://github.com/James-Yu/LaTeX-Workshop}

\end{hangparas}